\chapter{Mathematical Foundations and Notation}
\label{chap:foundations}

This chapter fixes the notation used throughout the notes and collects the mathematical concepts that are repeatedly required in probability and statistics.

\section{Sets and events}

A sample space is denoted by $\Omega$, while events are represented by subsets $A,B\subseteq\Omega$. Set union, intersection, complement, and set difference are denoted by $A\cup B$, $A\cap B$, $A^{c}$, and $A\setminus B$, respectively.

\section{Summation and products}

For a sequence $x_1,\ldots,x_n$, the finite sum and product are written as
\begin{align}
\sum_{i=1}^{n}x_i,
\qquad
\prod_{i=1}^{n}x_i.
\end{align}

Index notation should be kept explicit whenever the range is not unambiguous.

\section{Functions and transformations}

Random variables will be treated as measurable functions from a probability space into a measurable state space. Transformations of random variables therefore inherit the ordinary language of functions, inverse images, monotonicity, and change of variables.

\section{Linear algebra}

Vectors are represented by bold symbols such as $\bm{x}$, matrices by capital letters such as $X$, and transposition by $X^{\mathsf T}$. For regression and multivariate statistics, the Euclidean inner product is
\begin{align}
\langle \bm{x},\bm{y}\rangle = \bm{x}^{\mathsf T}\bm{y}.
\end{align}
